\section{五道门槛}

本报告将“悄悄布局尚未充分启动”拆成五个可验证条件，而不是凭K线观感选股：

\begin{enumerate}
  \item \textbf{价格位置}：行业或个股必须仍处一年价格区间中低位，或距离高点存在足够回撤。
  \item \textbf{资金持续性}：至少同时观察日度、周度和连续流入天数；单日净流入不能独立构成结论。
  \item \textbf{估值与盈利}：便宜必须有盈利稳定、边际改善或现金流支撑，否则归为低位无资金或价值陷阱。
  \item \textbf{拥挤度}：融资余额、成交放大和短期涨幅不能已经进入高拥挤区。
  \item \textbf{可执行估值}：核心标的必须有当前价格、2026E盈利分母、适配的估值方法、情景目标和失效条件。
\end{enumerate}

\section{数据口径}

\begin{exhibitbox}[Exhibit 2：数据源与适用边界]
\begin{tabularx}{\textwidth}{L{3.2cm}L{3.1cm}X}
\toprule
\textbf{数据} & \textbf{来源/质量} & \textbf{用于什么、不证明什么} \\
\midrule
行业历史价格 & 申万研究指数接口，L1 & 计算20/60/120日收益、52周位置；不证明资金身份 \\
行业估值分位 & 乐咕乐股申万快照，L2 & 判断PE/PB历史位置；不证明盈利即将反转 \\
日/周主力资金 & 证券时报数据宝结构化表，L2 & 判断交易资金方向；“主力”是成交单分类，不是实名机构 \\
连续流入个股 & 证券时报数据宝，L2 & 寻找价格尚未反应的资金孤岛；不等于推荐 \\
个股行情 & AStock新浪快照，L1 & 现价、振幅、成交额和日期 \\
公司财务 & 上市公司官方2026Q1报告，L1 & 已实现财务；不代替2026E预测 \\
盈利预测/目标 & 原始券商PDF及可审计完整转载，L2--L3 & 作为分母和外部锚；弱来源、内部推导目标均零权重 \\
\bottomrule
\end{tabularx}
\end{exhibitbox}

\section{数据降级与修复}

东方财富行业接口在本次采集中多次返回远端关闭连接，Baostock也出现网络接收错误。我们没有把失败解释为“没有资金数据”，而是切换到申万研究指数、证券时报数据宝和乐咕乐股，并保存原始HTML/JSON。申万历史接口最新记录停在7月9日，因此7月10日行业收盘点位采用7月9日官方收盘乘以数据宝披露的7月10日申万行业涨跌幅推算；报告中的7月10日\textbf{涨跌幅为原始披露值}，点位仅用于位置计算。

\section{分类定义}

\begin{itemize}
  \item \textbf{静默吸纳}：周度和日度资金均净流入，周涨幅受限，价格位置与估值未进入高位。
  \item \textbf{启动确认}：此前有资金承接，但单日价格已经显著上涨并放量。
  \item \textbf{单日反弹}：当日流入但周度仍净流出。
  \item \textbf{低位无资金}：价格低但资金持续流出。
  \item \textbf{资金孤岛}：一级行业整体未通过，但个别公司出现连续流入、低涨幅和基本面改善。
\end{itemize}
